\documentclass[handout_subfiles_root.tex]{subfiles}
\usepackage[rm2]{lecturenote}
\usetikzlibrary{decorations.pathreplacing}

\begin{document}
\HandoutTitle{11}{Sampling Theory}{September 29, 2026}
\maketitle

\HandoutMarkSection{1}{Sampling Theory: \texorpdfstring{$x(t)\longleftrightarrow x[n]$}{x(t) <-> x[n]}}

\noindent Consider the analog signal
\begin{equation*}
x(t) \longleftrightarrow X(j\Omega).
\end{equation*}
Sample uniformly at an interval $T$:
\begin{equation*}
x[n] = x(t)\big|_{t=nT}, \qquad n \text{ integer},
\qquad\qquad
x[n] \longleftrightarrow X(e^{j\omega}).
\end{equation*}

\begin{factbox}[title={Question}]
What is the relationship between $X(j\Omega)$ and $X(e^{j\omega})$?
\end{factbox}

\noindent\textbf{Recall:}
\begin{equation*}
x(t) = \frac{1}{2\pi}\int_{-\infty}^{\infty} X(j\Omega)\,e^{j\Omega t}\,d\Omega
\end{equation*}
\begin{align*}
\Rightarrow\quad x[n] &= \frac{1}{2\pi}\int_{-\infty}^{\infty} X(j\Omega)\,e^{j\Omega nT}\,d\Omega \\
&= \frac{1}{2\pi}\sum_{k=-\infty}^{\infty}\int_{-\pi/T}^{\pi/T} X\!\left(j\Big(\Omega+\frac{2\pi k}{T}\Big)\right) e^{j\left(\Omega+\frac{2\pi k}{T}\right)nT}\,d\Omega \\
&= \frac{1}{2\pi}\sum_{k=-\infty}^{\infty}\int_{-\pi/T}^{\pi/T} X\!\left(j\Big(\Omega+\frac{2\pi k}{T}\Big)\right) e^{j\Omega nT}\,d\Omega.
\end{align*}

\begin{notebox}[title=\HandoutNoteAddendumTitle]
The second line splits the real line into intervals of length $\tfrac{2\pi}{T}$ centered at $\tfrac{2\pi k}{T}$. The third line uses $e^{j\frac{2\pi k}{T}nT} = e^{j2\pi kn} = 1$ for integers $k,n$.
\end{notebox}

\noindent Change of variables:
\begin{equation*}
\omega = \Omega T \;\Longleftrightarrow\; \Omega = \frac{\omega}{T},
\qquad
\Omega = \pm\frac{\pi}{T} \;\Rightarrow\; \omega = \pm\pi,
\qquad
d\omega = T\,d\Omega.
\end{equation*}

\HandoutMarkSubsection{2}{The Sampled Spectrum}

\begin{align*}
x[n] &= \frac{1}{2\pi T}\sum_{k=-\infty}^{\infty}\int_{-\pi}^{\pi} X\!\left(j\frac{\omega}{T}+j\frac{2\pi k}{T}\right) e^{j\omega n}\,d\omega \\
&= \frac{1}{2\pi}\int_{-\pi}^{\pi} \underbrace{\left[\frac{1}{T}\sum_{k=-\infty}^{\infty} X\!\left(j\frac{\omega}{T}+j\frac{2\pi k}{T}\right)\right]}_{\HandoutUnderbraceLabel{X(e^{j\omega})}} e^{j\omega n}\,d\omega
\end{align*}
\vspace{0.3em}

\begin{factbox}[title={Sampled spectrum}]
\begin{equation*}
X(e^{j\omega}) = \frac{1}{T}\sum_{k=-\infty}^{\infty} X(j\Omega)\Big|_{\Omega=\frac{\omega}{T}+\frac{2\pi k}{T}}
\end{equation*}
\end{factbox}

\noindent Example: a triangular $X(j\Omega)$ of height $1$ supported on $[-\Omega_0,\Omega_0]$. The $k=0$ term is $\tfrac{1}{T}X\!\big(j\tfrac{\omega}{T}\big)$, a triangle of height $\tfrac1T$ on $[-\Omega_0T,\Omega_0T]$:

\begin{handouttikz}[x=0.9cm, y=0.9cm, font=\small]
  \begin{scope}
    \draw[->, ee483rule] (-3,0) -- (3,0) node[right, ee483ink] {$\Omega$};
    \draw[thick] (-1.2,0) -- (0,1.4) -- (1.2,0);
    \node[left] at (-2.4,1.2) {$X(j\Omega)$};
    \node[above] at (0,1.4) {$1$};
    \foreach \x/\l in {-1.2/{-\Omega_0}, 0/{0}, 1.2/{\Omega_0}} {\draw (\x,0.08) -- (\x,-0.08) node[below] {$\l$};}
  \end{scope}
  \begin{scope}[xshift=7cm]
    \draw[->, ee483rule] (-3,0) -- (3,0) node[right, ee483ink] {$\omega$};
    \draw[thick] (-1.6,0) -- (0,1.0) -- (1.6,0);
    \node[left] at (-2.2,1.2) {$k=0$:\ $\tfrac1T X\!\big(j\tfrac{\omega}{T}\big)$};
    \node[above] at (0,1.0) {$\tfrac1T$};
    \foreach \x/\l in {-1.6/{-\Omega_0T}, 0/{0}, 1.6/{\Omega_0T}} {\draw (\x,0.08) -- (\x,-0.08) node[below] {$\l$};}
  \end{scope}
\end{handouttikz}

\HandoutMarkParagraph{3}{The \texorpdfstring{$k=-1$}{k=-1} term and the sum over \texorpdfstring{$k$}{k}.}
\begin{equation*}
k=-1 \;\Rightarrow\; \frac{1}{T}X\!\left(j\frac{\omega}{T}-j\frac{2\pi}{T}\right)
= \frac{1}{T}X\!\left(j\Big(\Omega-\frac{2\pi}{T}\Big)\right)\bigg|_{\Omega=\frac{\omega}{T}},
\end{equation*}
i.e., the triangle moves to $\Omega=\tfrac{2\pi}{T}$ (support $\big[\tfrac{2\pi}{T}-\Omega_0,\tfrac{2\pi}{T}+\Omega_0\big]$), which is $\omega=2\pi$ (support $[2\pi-\Omega_0T,\,2\pi+\Omega_0T]$). Each $k$ contributes one copy centered at $\omega=-2\pi k$:

\begin{handouttikz}[x=0.62cm, y=0.8cm, font=\small]
  \foreach \row/\w in {0/0.75, -2.4/1.45} {
    \draw[->, ee483rule] (-11.2,\row) -- (11.2,\row) node[right, ee483ink] {$\omega$};
    \foreach \c in {-8,-4,0,4,8} {\draw[thick] (\c-2*\w,\row) -- (\c,\row+1.2) -- (\c+2*\w,\row);}
    \foreach \c/\l in {-8/{-4\pi}, -4/{-2\pi}, 0/{0}, 4/{2\pi}, 8/{4\pi}} {\draw (\c,\row+0.08) -- (\c,\row-0.08) node[below] {$\l$};}
  }
  \foreach \c/\k in {-8/{+2}, -4/{+1}, 0/{0}, 4/{-1}, 8/{-2}} {\node[above] at (\c,1.25) {$k=\k$};}
  \node[right] at (0.2,-1.1) {$\tfrac1T$};
\end{handouttikz}
{\centering\scriptsize Top: copies do not overlap. Bottom: wider $X(j\Omega)$ (or larger $T$), so neighboring copies overlap.\par}

\HandoutMarkSection{4}{Summary: CTFT, Sampling, and DTFT}

\begin{handouttikz}[x=1cm, y=1cm, font=\small]
  \node (xt) at (0,0) {$x(t)$};
  \node (Xj) at (3.2,0) {$X(j\Omega)$};
  \node (xn) at (0,-1.8) {$x[n] = x(t)\big|_{t=nT}$};
  \node (Xe) at (4.6,-1.3) {$X(e^{j\omega})$};
  \draw[<->, thick] (xt) -- node[above] {CTFT} (Xj);
  \draw[->, thick] (xt) -- node[right] {sampling} (xn);
  \draw[<->, thick] (xn.north east) ++(0,-0.15) -- node[above] {DTFT} (Xe.west);
\end{handouttikz}

\begin{equation*}
X(e^{j\omega}) = \underbrace{\frac{1}{T}}_{\HandoutUnderbraceLabel{\text{rescaling}}}\ \underbrace{\sum_{k=-\infty}^{\infty}}_{\HandoutUnderbraceLabel{\text{periodization}}} X(j\Omega)\Big|_{\underbrace{\scriptstyle\Omega=\frac{\omega}{T}+\frac{2\pi k}{T}}_{\HandoutUnderbraceLabel{\text{stretching of the axis}}}}
\end{equation*}
\vspace{0.3em}

\noindent Two cases of interest:
\begin{itemize}
  \item $X(j\Omega) = 0$ whenever $|\Omega| \ge \frac{\pi}{T}$ \quad (bandlimited signal with bandwidth $B \le \frac{2\pi}{T}$).
  \item $X(j\Omega) \ne 0$ for some $\Omega \ge \frac{\pi}{T}$ or $\Omega \le -\frac{\pi}{T}$.
\end{itemize}

\begin{handouttikz}[x=0.9cm, y=0.9cm, font=\small]
  \begin{scope}
    \draw[->, ee483rule] (-3,0) -- (3,0) node[right, ee483ink] {$\Omega$};
    \draw[thick] (-1.2,0) -- (0,1.3) -- (1.2,0);
    \node[above] at (0,1.3) {$1$};
    \foreach \x/\l in {-1.2/{-\frac{\pi}{T}}, 0/{0}, 1.2/{\frac{\pi}{T}}} {\draw (\x,0.08) -- (\x,-0.08) node[below] {$\l$};}
    \draw[decorate, decoration={brace, mirror, amplitude=4pt}] (-1.2,-0.75) -- node[below=4pt] {$B$} (1.2,-0.75);
    \node at (0,-1.9) {bandlimited};
  \end{scope}
  \begin{scope}[xshift=7cm]
    \draw[->, ee483rule] (-3,0) -- (3,0) node[right, ee483ink] {$\Omega$};
    \draw[thick] (-2.4,0) -- (0,1.3) -- (2.4,0);
    \foreach \x/\l in {-1.2/{-\frac{\pi}{T}}, 0/{0}, 1.2/{\frac{\pi}{T}}} {\draw (\x,0.08) -- (\x,-0.08) node[below] {$\l$};}
    \node at (0,-1.9) {not bandlimited to $\frac{\pi}{T}$};
  \end{scope}
\end{handouttikz}

\HandoutMarkSubsection{5}{The Bandlimited Case}

\noindent Consider
\begin{equation*}
\frac{1}{T}\sum_{k=-\infty}^{\infty} X(j\Omega)\Big|_{\Omega=\tilde\Omega+\frac{2\pi k}{T}}.
\end{equation*}

\begin{handouttikz}[x=0.62cm, y=0.8cm, font=\small]
  \foreach \row/\s/\v in {0/{\tilde\Omega}/T, -2.6/{\omega}/1} {
    \draw[->, ee483rule] (-10.5,\row) -- (10.5,\row) node[right, ee483ink] {$\s$};
    \foreach \c in {-6,0,6} {\draw[thick] (\c-3,\row) -- (\c,\row+1.3) -- (\c+3,\row);}
    \draw[thick] (-10.2,\row+0.4) -- (-9,\row); \draw[thick] (9,\row) -- (10.2,\row+0.4);
  }
  \foreach \c/\l in {-9/{-\frac{3\pi}{T}}, -6/{-\frac{2\pi}{T}}, -3/{-\frac{\pi}{T}}, 0/{0}, 3/{\frac{\pi}{T}}, 6/{\frac{2\pi}{T}}, 9/{\frac{3\pi}{T}}} {\draw (\c,0.08) -- (\c,-0.08) node[below] {$\l$};}
  \foreach \c/\l in {-9/{-3\pi}, -6/{-2\pi}, -3/{-\pi}, 0/{0}, 3/{\pi}, 6/{2\pi}, 9/{3\pi}} {\draw (\c,-2.52) -- (\c,-2.68) node[below] {$\l$};}
  \node[above] at (0,1.3) {$\tfrac1T$};
  \node[above] at (-6,1.3) {$k=+1$}; \node[above] at (6,1.3) {$k=-1$}; \node[above right] at (0.2,1.0) {$k=0$};
  \node[left] at (-10.5,-1.6) {$\omega=\tilde\Omega T$};
  \draw[decorate, decoration={brace, mirror, amplitude=4pt}] (-3,-3.5) -- (3,-3.5);
\end{handouttikz}

\noindent On $-\pi\le\omega\le\pi$, the \textbf{shape of the original CTFT $X(j\Omega)$ is preserved} (and $x(t)$ can be perfectly recovered from $x[n]$ if we follow appropriate steps).

\Needspace{16\baselineskip}
\HandoutMarkSubsection{6}{The Other Case (Aliasing) and the Sampling Theorem}

\begin{handouttikz}[x=0.62cm, y=0.8cm, font=\small]
  % shifted copies: each wider than 2pi/T, so neighbors overlap
  \begin{scope}
    \clip (-10.5,-0.2) rectangle (10.5,1.7);
    \foreach \c in {-16,-8,0,8,16} {\draw[thick, ee483rule] (\c-5.5,0) -- (\c,1.4) -- (\c+5.5,0);}
  \end{scope}
  \draw[->, ee483rule] (-10.5,0) -- (10.5,0) node[right, ee483ink] {$\tilde\Omega$};
  \foreach \c/\l in {-8/{-\frac{2\pi}{T}}, -4/{-\frac{\pi}{T}}, 0/{0}, 4/{\frac{\pi}{T}}, 8/{\frac{2\pi}{T}}} {\draw (\c,0.08) -- (\c,-0.08) node[below] {$\l$};}
  % the resulting sum
  \begin{scope}[yshift=-2.8cm]
    \begin{scope}
      \clip (-10.5,-0.2) rectangle (10.5,2.2);
      \draw[thick] plot[domain=-10.5:10.5, samples=211]
        (\x, {1.4*(max(0,1-abs(\x+16)/5.5)+max(0,1-abs(\x+8)/5.5)+max(0,1-abs(\x)/5.5)+max(0,1-abs(\x-8)/5.5)+max(0,1-abs(\x-16)/5.5))});
    \end{scope}
    \draw[->, ee483rule] (-10.5,0) -- (10.5,0) node[right, ee483ink] {$\tilde\Omega$};
    \foreach \c/\l in {-8/{-\frac{2\pi}{T}}, -4/{-\frac{\pi}{T}}, 0/{0}, 4/{\frac{\pi}{T}}, 8/{\frac{2\pi}{T}}} {\draw (\c,0.08) -- (\c,-0.08) node[below] {$\l$};}
  \end{scope}
  \foreach \x in {-5.5,-2.5,2.5,5.5} {\draw[densely dashed, ee483rule] (\x,0) -- (\x,-2.25);}
  % same axis in omega
  \draw[->, ee483rule] (-10.5,-5.2) -- (10.5,-5.2) node[right, ee483ink] {$\omega$};
  \foreach \c/\l in {-8/{-2\pi}, -4/{-\pi}, 0/{0}, 4/{\pi}, 8/{2\pi}} {\draw (\c,-5.12) -- (\c,-5.28) node[below] {$\l$};}
\end{handouttikz}
{\centering\scriptsize Overlapping copies add up, so the sum no longer looks like one copy of $X(j\Omega)$.\par}

\noindent The \textbf{CTFT shape is not preserved} (EE~592).

\begin{defbox}[title={Sampling theorem}]
A signal can be recovered from its samples provided that the signal satisfies
\begin{equation*}
X(j\Omega) = 0 \quad \text{for } \forall\,|\Omega| \ge \frac{\pi}{T},
\end{equation*}
where $T$ is the sampling interval.
\begin{itemize}
  \item Nyquist \qquad -- Shannon \qquad -- Whittaker \qquad -- Kotelnikov
\end{itemize}
\end{defbox}

\HandoutMarkSection{7}{Sampling Rate and the Nyquist Criterion}

\noindent Converting from $\Omega$ (radians/sec) to $f$ (Hz $=1/\text{sec}$):
\begin{equation*}
f = \frac{\Omega}{2\pi\ \text{radians}}.
\end{equation*}
Largest signal frequency:
\begin{equation*}
|f_{\max}| < \frac{(\pi/T)}{2\pi} = \frac{1}{2T}.
\end{equation*}
Let $f_s \triangleq \frac{1}{T}$ be the \textbf{sampling rate} (sampling frequency).

\begin{factbox}[title={Nyquist criterion (Nyquist condition)}]
For no aliasing,
\begin{equation*}
f_s > 2|f_{\max}|.
\end{equation*}
\begin{itemize}
  \item Frequency $2|f_{\max}|$: \textbf{Nyquist rate}.
  \item Frequency $f_s/2$: \textbf{Nyquist frequency}.
\end{itemize}
\end{factbox}

\begin{handouttikz}[x=1cm, y=1cm, font=\small]
  \node (in) at (0,0) {$x(t)$};
  \node[draw, minimum width=1.4cm, minimum height=0.8cm, align=center] (ad) at (2.6,0) {ideal\\A/D};
  \node[draw, minimum width=1.4cm, minimum height=0.8cm, align=center] (da) at (8,0) {ideal\\D/A};
  \node (out) at (10.6,0) {$x(t)$};
  \draw[->, thick] (in) -- (ad);
  \draw[->, thick] (ad) -- node[above] {$x[n]=x(t)\big|_{t=nT}$} (da);
  \draw[->, thick] (da) -- (out);
  \draw[rounded corners, thick] (1.5,-1.3) rectangle (9.1,1.1);
  \node[above left] at (9.1,1.1) {digital ``box''};
  \node[above] at (1.2,1.1) {analog signal};
  \node[align=center, font=\scriptsize] at (2.6,-1.0) {sampling at interval $T$};
\end{handouttikz}

\noindent This works if $x(t)$ is appropriately bandlimited.

\HandoutMarkSubsection{8}{Discrete-Time Processing of Analog Signals}

\begin{handouttikz}[x=1cm, y=1cm, font=\small]
  \node (in) at (0,0) {$x(t)$};
  \node[draw, minimum height=0.8cm, align=center] (ad) at (2.1,0) {ideal\\A/D};
  \node[draw, minimum height=0.8cm, align=center] (dt) at (5.4,0) {discrete-time\\processing};
  \node[draw, minimum height=0.8cm, align=center] (da) at (8.7,0) {ideal\\D/A};
  \node (out) at (10.7,0) {$y(t)$};
  \draw[->, thick] (in) -- (ad);
  \draw[->, thick] (ad) -- node[above] {$x[n]$} (dt);
  \draw[->, thick] (dt) -- (da);
  \draw[->, thick] (da) -- (out);
  \draw[rounded corners, thick] (1.3,-0.9) rectangle (9.5,0.9);
\end{handouttikz}

\noindent Equivalent to analog processing (but is much easier).

\noindent\textbf{What if $x(t)$ is not bandlimited?}

\begin{handouttikz}[x=1cm, y=1cm, font=\small]
  \node (in) at (0,0) {$x(t)$};
  \node[draw, minimum height=0.8cm] (h) at (2.2,0) {$H(j\Omega)$};
  \node[draw, minimum height=0.8cm, align=center] (p) at (7.9,0) {processing\\from before};
  \node (out) at (10.4,0) {$\tilde x(t)$};
  \draw[->, thick] (in) -- (h);
  \draw[->, thick] (h) -- node[above] {$\tilde x(t)$} (p);
  \draw[->, thick] (p) -- (out);
  \node[below, align=center] at (2.2,-0.5) {ideal LPF\\(analog antialiasing filter)};
  \node[below, align=center] at (5.3,-0.1) {$\tilde x(t)$ is truly\\bandlimited};
\end{handouttikz}

\begin{examplebox}[title={Example}]
\begin{align*}
g_1(t) &= \cos(6\pi t) &&\longleftrightarrow\quad G_1(j\Omega) = \pi\big[\delta(\Omega-6\pi)+\delta(\Omega+6\pi)\big], \\
g_2(t) &= \cos(14\pi t) &&\longleftrightarrow\quad G_2(j\Omega) = \pi\big[\delta(\Omega-14\pi)+\delta(\Omega+14\pi)\big], \\
g_3(t) &= \cos(26\pi t) &&\longleftrightarrow\quad G_3(j\Omega) = \pi\big[\delta(\Omega-26\pi)+\delta(\Omega+26\pi)\big].
\end{align*}
These are $3$\,Hz, $7$\,Hz, and $13$\,Hz signals. Sample at interval $T=0.1$\,sec:
\begin{equation*}
\text{sampling freq} = \frac{1}{T} = 10\text{ Hz} = 20\pi \text{ rad/sec}.
\end{equation*}

\HandoutMarkLeft{9}\noindent Signal content above $5$\,Hz will be aliased.

\begin{center}
\begin{tikzpicture}[x=0.36cm, y=0.7cm, font=\small]
  % G1
  \draw[->, ee483rule] (-11,0) -- (11,0) node[right, ee483ink] {$\Omega$};
  \foreach \x in {-6,6} {\draw[->, thick] (\x,0) -- (\x,1.2);}
  \foreach \x/\l in {-10/{-10\pi}, -6/{-6\pi}, 0/{0}, 6/{6\pi}, 10/{10\pi}} {\draw (\x,0.08) -- (\x,-0.08) node[below] {$\l$};}
  \node[left] at (-11,0.9) {$G_1(j\Omega)$};
  \draw[->, ee483rule] (-11,-2.4) -- (11,-2.4) node[right, ee483ink] {$\omega$};
  \foreach \x in {-6,6} {\draw[->, thick] (\x,-2.4) -- (\x,-1.2);}
  \foreach \x/\l in {-10/{-\pi}, -6/{-0.6\pi}, 0/{0}, 6/{0.6\pi}, 10/{\pi}} {\draw (\x,-2.32) -- (\x,-2.48) node[below] {$\l$};}
\end{tikzpicture}

\vspace{0.6em}
\begin{tikzpicture}[x=0.3cm, y=0.7cm, font=\small]
  % G2
  \draw[->, ee483rule] (-16,0) -- (16,0) node[right, ee483ink] {$\Omega$};
  \foreach \x in {-14,14} {\draw[->, thick] (\x,0) -- (\x,1.2);}
  \foreach \x/\l in {-14/{-14\pi}, -10/{-10\pi}, 0/{0}, 10/{10\pi}, 14/{14\pi}} {\draw (\x,0.08) -- (\x,-0.08) node[below] {$\l$};}
  \node[left] at (-16,0.9) {$G_2(j\Omega)$};
  \draw[->, ee483rule] (-16,-2.6) -- (16,-2.6) node[right, ee483ink] {$\omega$};
  \foreach \x in {-14,-6,6,14} {\draw[->, thick] (\x,-2.6) -- (\x,-1.4);}
  \foreach \x/\l in {-14/{-1.4\pi}, -10/{-\pi}, -6/{-0.6\pi}, 0/{0}, 6/{0.6\pi}, 10/{\pi}, 14/{1.4\pi}} {\draw (\x,-2.52) -- (\x,-2.68) node[below] {$\l$};}
  \draw[->, thin, ee483link] (14,-1.3) to[bend right=25] (-6,-1.3);
  \draw[->, thin, ee483link] (-14,-1.3) to[bend left=25] (6,-1.3);
\end{tikzpicture}
\end{center}

\noindent For $G_3(j\Omega)$: $\Omega = \pm 26\pi \;\Rightarrow\; \omega = \pm 2.6\pi \;\Rightarrow\; \pm 0.6\pi$.

\noindent\textbf{We can't tell the difference between any of these signals!} \quad $g_1[n] = g_2[n] = g_3[n]$.
\end{examplebox}

\begin{notebox}[title=\HandoutNoteAddendumTitle]
Directly in time: $g_2[n] = \cos(1.4\pi n) = \cos(1.4\pi n - 2\pi n) = \cos(-0.6\pi n) = \cos(0.6\pi n) = g_1[n]$, and likewise $g_3[n] = \cos(2.6\pi n) = \cos(2.6\pi n - 2\pi n) = \cos(0.6\pi n)$. Only $g_1$ ($3$\,Hz $<5$\,Hz) satisfies the Nyquist criterion, so $3$\,Hz is the frequency an ideal D/A would reconstruct for all three.
\end{notebox}

\section*{Readings}
\begin{itemize}
  \item Sanjit K. Mitra, \textit{Digital Signal Processing: A Computer-Based Approach}, 4th ed., Secs.\ 2.5, 3.8, 3.9
  \item Monson H. Hayes, \textit{Schaum's Outline of Digital Signal Processing}, 2nd ed., Ch.\ 3
\end{itemize}

\HandoutAppendixListStart
  \HandoutAppendixListItem{app:scan11}{A.\ Handwritten Lecture Note Scan (September 29, 2026)}
\HandoutAppendixListEnd

\HandoutAppendixSection{app:scan11}{Handwritten Lecture Note Scan}
\HandoutIncludeHandoutScan{lecture11_0929.pdf}

\end{document}
